\documentclass[a4paper,10pt]{article}
\usepackage[left=0cm,top=3cm,right=2.5cm,bottom=4cm]{geometry} 
\usepackage[utf8]{inputenc}
\usepackage[activeacute,spanish]{babel}
\usepackage{multirow}	% Para poder unir filas en las tablas
\usepackage{colortbl}	% Para colorear tablas
\usepackage{latexsym,amssymb,oldlfont,amsmath}
\usepackage{graphicx,graphics}
\usepackage{multicol}
\usepackage{color}
\usepackage{listings}
\usepackage{lastpage}
\usepackage{enumitem}



%\usepackage{fontspec}
%\setmainfont{Times New Roman}
\newcommand{\R}{{\mathbb R}}
\newcommand{\si}{{\mathcal S}^n}
\renewcommand{\theequation}{\thesection.\arabic{equation}} % Cambia formato salida numeracion equation
\makeatletter\@addtoreset {equation}{section} \makeatother % Tambien se usa
\renewcommand{\baselinestretch}{1.3}
\usepackage{amssymb}
\usepackage{hyperref}
\definecolor{midnightblue}{rgb}{0.1,0.1,0.44}


\newcommand{\dis}{\displaystyle}
\renewcommand{\theenumii}{\arabic{enumii}}

\def\Z{\mathbb{Z}}
\def\N{\mathbb{N}}
\def\R{\mathbb{R}}
\def\M{\mathcal{M}_n(\mathbb{R})}
\def\sin{\operatorname{sen}}
\def\e{\operatorname{e}}
\DeclareMathOperator{\arcsec}{arcsec}
\DeclareMathOperator{\arccot}{arccot}
\DeclareMathOperator{\arccsc}{arccsc}



\usepackage{etoolbox,fancyhdr,xcolor}
\fancyhf{} % sets both header and footer to nothing
\renewcommand{\headrulewidth}{0pt}
%\renewcommand{\footrulewidth}{0.5pt}
\pagestyle{fancy}
\headheight=60pt %para cambiar el tamaño del encabezadox
\fancyhead[L] %la "L" indica a la izquierda
{
\begin{minipage}{6cm}
\includegraphics[width=0.9\textwidth]{logo_udp.eps}
\end{minipage}
}

%\fancyhead[C]
%{
%\textbf{Cálculo II - Guía 1}
%}


%\fancyhead[R] %la "L" indica a la izquierda
%{

%}
%\begin{minipage}{4cm}
%\small %tamaño de letra pequeño
%{
%\textsf
%{
%Universidad de Chile\\
%...}}
%\end{minipage}}



%\fancyfoot[C]{}
%{
%\begin{minipage}{10cm}
%\color{midnightblue}
%\begin{center}Casa Matriz: Monjitas 527 Piso 15 - Santiago de Chile\\
%Fono (02)24367255 \end{center}
%\end{minipage}
%}
\fancyfoot[C]{}
%Página \thepage\ de \pageref{LastPage}}% Páginas Pares numeradas a la derecha
%\fancyfoot[LO]{\thepage} % Páginas Impares numeradas a la izquierda

\title{C2_Ay4_pauta}

\begin{document}

\begin{center}
Cálculo II\\
Ayudantía Nº4 - Pauta\\
Primer Semestre 2017
\end{center}



\begin{enumerate}
\item Justifique la veracidad del siguiente enunciado:

Una empresa tiene como función de ingreso marginal a $I'(x)=35-4x$ y función de costo marginal a $C'(x) = 5-2x$, donde $x$ es la cantidad de artículos vendidos, y las unidades monetarias están en millones de pesos [MM\$]. Suponga que cuando no se venden artículos el ingreso total es cero y el costo total es $25$[MM\$] (por conceptos de gastos fijos). Recordando que la utilidad de una empresa es la diferencia entre los ingresos totales y los costos totales, entonces la cantidad de artículos a vender que maximizan las utilidades son $225$ artículos. \\

Datos:
\begin{align*}
 \dis  I'(x)=35-4x \\
 \dis  C'(x) = 5-2x \\
 \dis I(0) = 0\\
 \dis C(0) = 25
\end{align*}
Calculando el ingreso:
\begin{align*}
 \dis  I(x)  &=  \dis  \int (35-4x) \, dx  \\
 \dis  I(x) &=  \dis  35x - 2x^2 + C  \\
 \dis  I(0) &=  C  \\
 \dis  C = 0
\end{align*}
Función ingreso total:
\begin{align*}
 \dis  I(x)  &=  \dis  35x - 2x^2
\end{align*}

Calculando costos:
\begin{align*}
 \dis  C(x)  &=  \dis  \int (5-2x) \, dx  \\
 \dis  C(x) &=  \dis  5x - x^2 + C  \\
 \dis  C(0) &=  C  \\
 \dis  C = 25
\end{align*}
Función costo total:
\begin{align*}
 \dis  C(x)  &=  \dis  5x - x^2 + 25
\end{align*}
\dotfill\textbf{[1 punto]}

Utilidad:
\begin{align*}
 \dis  U(x)  &=  \dis  I(x)-C(x)\\
 \dis  U(x)  &=  \dis  35x - 2x^2-(5x - x^2 + 25)\\
 \dis  U(x)  &=  \dis  30x - x^2 - 25\\
\end{align*}
Buscando la utilidad máxima:
Utilidad:
\begin{align*}
 \dis  U'(x)  &=  30 - 6x\\
  \dis  30 - 2x  &= 0 \\
  \dis  x  &= 15 \\
\end{align*}
Las utilidades se maximizan para $x= 15$. En efecto, $U''(15)=-2$ lo cual es menor que cero.
Así, vendiendo 15 articulos se maximizan las utilidades, no 225 artículos. Por lo tanto la afirmación es \underline{FALSA}.

\dotfill\textbf{[1 punto]}

\newpage

\item Considere la función $f(x) = |\sin(x)|$ definida para $x \in [0,2\pi]$, y la partición $\displaystyle P = \left\{ 0, \frac{\pi}{4} , \frac{2\pi}{3} , \pi , \frac{3\pi}{2} , \frac{11\pi}{6} , 2\pi \right\}$. Calcule la suma superior e inferior de $f$ sobre $P$. \\

\begin{align*}
 \dis  S_{inf}  &=  f(0)*(\frac{\pi}{4}-0) + f(\frac{\pi}{4})*(\frac{2\pi}{3}-\frac{\pi}{4}) + f(\frac{2\pi}{3})*(\pi-\frac{2\pi}{3}) + f(\pi)*(\frac{3\pi}{2}-\pi) + f(\frac{3\pi}{2})*(\frac{11\pi}{6}-\frac{3\pi}{2}) + f(\frac{11\pi}{6})*(2\pi-\frac{11\pi}{6})\\
  &=  0*(\frac{\pi}{4}) + \frac{\sqrt{2}}{2}*(\frac{5\pi}{12}) + \frac{\sqrt{3}}{2}*(\frac{\pi}{3}) + 0*(\frac{\pi}{2}) + (1)*(\frac{\pi}{3}) + (\frac{1}{2})*(\frac{\pi}{6})\\
  &=   \frac{5\sqrt{2}\pi}{24} + \frac{\sqrt{3}\pi}{6} + \frac{\pi}{3} + \frac{\pi}{12}
\end{align*}
\dotfill\textbf{[1 punto]}

\begin{align*}
 \dis  S_{sup}  &=  f(\frac{\pi}{4})*(\frac{\pi}{4}-0) + f(\frac{2\pi}{3})*(\frac{2\pi}{3}-\frac{\pi}{4}) + f(\pi)*(\pi-\frac{2\pi}{3}) + f(\frac{3\pi}{2})*(\frac{3\pi}{2}-\pi) + f(\frac{11\pi}{6})*(\frac{11\pi}{6}-\frac{3\pi}{2}) + f(2\pi)*(2\pi-\frac{11\pi}{6})\\
  &=  \frac{\sqrt{2}}{2}*(\frac{\pi}{4}) + \frac{\sqrt{3}}{2}*(\frac{5\pi}{12}) + 0*(\frac{\pi}{3}) + (1)*(\frac{\pi}{2}) + (\frac{1}{2})*(\frac{\pi}{3}) + 0*(\frac{\pi}{6})\\
  &=   \frac{\sqrt{2}\pi}{8} + \frac{5\sqrt{3}\pi}{24} + \frac{\pi}{2} + \frac{\pi}{6}
\end{align*}
\dotfill\textbf{[1 punto]}

\newpage

\item Considere la función $f(x)= [\sqrt{x}]$, donde $[x]$ denota la parte entera de $x$.
\begin{enumerate}[label=\emph{\alph*)}]
\item Grafique la función para $x \in [0,16]$

\begin{figure}[ht]
    \centering
    \includegraphics[width=10cm]{c2}
    \label{3a}
\end{figure}
\dotfill\textbf{[1 punto]}

\item Obtenga una partición regular P del intervalo $[0,16]$ en 8 subintervalos.

\begin{align*}
 \dis  \Delta x &= \frac{(16-0)}{8}\\
 \dis  \Delta x &= 2
\end{align*}

Con lo que encontramos la partición $\displaystyle P = \left\{ 0, 2, 4, 6, 8, 10, 12, 14, 16 \right\}$

\dotfill\textbf{[1 punto]}

\item  Calcule la suma superior e inferior de $f$ acotada a la partición P.

\begin{align*}
 \dis  S_{inferior}  &=  f(0)*(2) + f(2)*(2) + f(4)*(2) + f(6)*(2) + f(8)*(2) + f(10)*(2) + f(12)*(2) + f(14)*(2)\\
  &= 0*(2) + 1*(2) + 2*(2) + 2*(2) + 2*(2) + 3*(2) + 3*(2) + 3*(2)\\
  &= 32
\end{align*}

\dotfill\textbf{[1 punto]}

\begin{align*}
 \dis  S_{superior}  &=  f(2)*(2) + f(4)*(2) + f(6)*(2) + f(8)*(2) + f(10)*(2) + f(12)*(2) + f(14)*(2) + f(16)*(2)\\
  &= 1*(2) + 2*(2) + 2*(2) + 2*(2) + 3*(2) + 3*(2) + 3*(2) + 4*(2)\\
  &= 40
\end{align*}

\dotfill\textbf{[1 punto]}

\end{enumerate}


\newpage
\item Mediante sumas de Riemann, calcule el área bajo la curva $f(x) = 9-x^2$ en el intervalo $[-3,3]$\\

Utilizando las sumas de Riemann:

\begin{align*}
  \dis  \displaystyle \int_{-3}^3 (9-x^2) \, dx  &=  \displaystyle\lim_{n \to \infty } \frac{3-(-3)}{n} \cdot \sum_{i=1}^n \left(9 - (-3+i\frac{6}{n})^2\right) \hspace{2mm} \textbf{[1 punto]} \\
    \dis    &=  \displaystyle\lim_{n \to \infty } \frac{6}{n}\sum_{i=1}^n \left(9 - \left(9-\frac{36i}{n}+\frac{36i^2}{n^2}\right)\right)\\
    \dis    &=  \displaystyle\lim_{n \to \infty } \frac{6}{n}\sum_{i=1}^n 36\left(\frac{i}{n}+\frac{i^2}{n^2}\right)\\
    \dis    &=  \displaystyle\lim_{n \to \infty } \frac{216}{n}\left( \frac{1}{n} \sum_{i=1}^n i - \frac{1}{n^2} \sum_{i=1}^n i^2\right) \\
\dis    &=  \displaystyle\lim_{n \to \infty } \frac{216}{n}\left( \frac{1}{n} \frac{n(n+1)}{2} - \frac{1}{n^2} \frac{n(n+1)(2n+1)}{6}\right) \hspace{2mm} \textbf{[1 punto]} \\
\dis    &=  \displaystyle\lim_{n \to \infty } 216\left( \frac{n+1}{2n} - \frac{2n^2+3n+1}{6n^2}\right)\\
\dis    &=   216\left( \frac{1}{2} - \frac{1}{3}\right)\\
\dis    &=   36 \hspace{2mm} \textbf{[1 punto]}
\end{align*}


\newpage
\item Mediante sumas de Riemann, calcule el área bajo la curva $f(x) = x^3$ en el intervalo $[0,1/2]$

\begin{align*}
  \dis  \displaystyle \int_{0}^{\frac{1}{2}} x^3 \, dx  &=  \displaystyle\lim_{n \to \infty } \frac{1}{2n} \cdot \sum_{i=1}^n \left( \frac{i}{2n} \right)^3 \hspace{2mm} \textbf{[1 punto]} \\
    \dis    &=  \displaystyle\lim_{n \to \infty } \frac{1}{2n} \cdot \frac{1}{8n^3}\sum_{i=1}^n i^3\\
\dis    &=  \displaystyle\lim_{n \to \infty } \frac{1}{16n^4} \left(\frac{n(n+1)}{2}\right)^2 \hspace{2mm} \textbf{[1 punto]} \\
\dis    &=  \frac{1}{64} \hspace{2mm} \textbf{[1 punto]}
\end{align*}


\end{enumerate}




\vspace{10mm}




\end{document}
